\textbf{COLUNI 2023 -- Questão 33}\par\medskip Leia o poema a seguir.

\begin{center}
\textbf{O Mundo do Menino Impossível}\\
\textbf{Jorge de Lima}
\end{center}

\noindent
\begin{minipage}[t]{0.32\linewidth}
Fim da tarde, boquinha da noite\\
Com as primeiras estrelas\\
E os derradeiros sinos.\\[0.6em]
Entre as estrelas e lá detrás da igreja,\\
Surge a lua cheia\\
Para chorar com os poetas.\\[0.6em]
E vão dormir as duas coisas novas desse mundo:\\
O sol e os meninos.\\[0.6em]
Mas ainda vela\\
O menino impossível\\
Aí do lado\\
Enquanto todas as crianças mansas\\
Dormem\\
Acalentadas\\
Por Mãe-negra noite.\\[0.6em]
O menino impossível\\
Que destruiu\\
Os brinquedos perfeitos\\
Que os vovós lhe deram:\\
O urso de Nüremberg,\\
O velho barbado iugoslavo,\\
As \textit{poupées de Paris aux}\\
\textit{Cheveux crêpés}
\end{minipage}\hfill
\begin{minipage}[t]{0.32\linewidth}
O carrinho português\\
Feito de folha-de-flandres,\\
A caixinha de música tcheco-eslovaca,\\
O polichinelo italiano\\
\textit{Made in England},\\
O trem de ferro de U.S.A\\
E o macaco brasileiro\\
De Buenos-Aires\\
\textit{Moviendo la cola y la cabeza.}\\[0.6em]
O menino impossível\\
Que destruiu até\\
Os soldados de chumbo de Moscou\\
E furou os olhos de um Papá Noel,\\
Brinca com sabugos de milho,\\
Caixas vazias,\\
Tacos de pau,\\
Pedrinhas brancas do rio...\\[0.6em]
“Faz de conta que os sabugos\\
São bois...”\\
“Faz de conta...”\\
“Faz de conta...”\\[0.6em]
E os sabugos de milho\\
Mugem como bois de verdade...\\[0.6em]
E os tacos que deveriam ser\\
Soldadinhos de chumbo são\\
cangaceiros de chapéus de couro...
\end{minipage}\hfill
\begin{minipage}[t]{0.32\linewidth}
E as pedrinhas balem!\\
Coitadinhas das ovelhas mansas\\
Longe das mães\\
Presas nos currais de papelão!\\[0.6em]
É boquinha da noite\\
No mundo que o menino impossível\\
Povoou sozinho!\\[0.6em]
A mamãe cochila.\\
O papai cabeceia.\\
O relógio badala.\\[0.6em]
E vem descendo\\
Uma noite encantada\\
Da lâmpada que expira\\
Lentamente\\
Na parede da sala...\\[0.6em]
O menino poisa a testa\\
E sonha dentro da noite quieta\\
Da lâmpada apagada\\
Com o mundo maravilhoso\\
Que ele tirou do nada...\\[0.6em]
Xô! Xô! Pavão!\\
Sai de cima do telhado!\\
Deixa o menino dormir\\
Seu soninho sossegado!
\end{minipage}

\par\medskip NUNES, Cassiano; BRITO, Mário da Silva. \textbf{Poesia Brasileira}. Editora Clubinho Literário, 2021. p. 145-147.

O mundo moderno encontra-se de tal modo integrado em redes de diversos tipos, num intenso processo de globalização, que em vários âmbitos encontram-se situações semelhantes a do menino impossível diante dos brinquedos ganhados dos avós...

Sobre os diversos brinquedos do menino e sua origem, é INCORRETO afirmar:
\begin{enumerate}
\item São originários de países pertencentes a diferentes continentes.
\item São oriundos de atividades industriais que processam diferentes matérias primas.
\item Dependem de um intrincado mecanismo de transporte para chegar às mãos do menino.
\item Comprovam que ele adota um padrão de vida consumista típico de países desenvolvidos.
\end{enumerate}
